% interaction/fiat_shamir.tex — interaction-native Fiat-Shamir

\section{Fiat--Shamir in the Dependent Interaction Core}\label{sec:interaction-fiat-shamir}

The Fiat--Shamir transform is a natural test case for the interaction-native
framework.  It is simple enough to explain concretely, but subtle enough to
reveal which parts of the framework are merely convenient and which parts are
mathematically necessary.

At a high level, the transform has the same familiar shape as in the flat
oracle-reduction setting: replace verifier challenges by deterministic values
derived from the transcript prefix, let the prover compute all of its messages
without further interaction, and let the verifier reconstruct the interactive
transcript before checking it.  The difference is that, in the dependent core,
the \emph{rest of the protocol} may depend on those challenge values.  So the
transform must do more than merely recover the right list of challenges: it
must recover the right \emph{subtree} of the interaction specification at each
receiver node.

\subsection{Replay oracles and messages-only proofs}

The key input to the basic transform is a deterministic \emph{replay oracle}:
given the current transcript prefix, it returns exactly the verifier challenge
that Fiat--Shamir would have produced there.

\begin{definition}[ReplayOracle]
  \label{int:replay-oracle}
  For a role-decorated interaction $(\mathit{spec}, \mathit{roles})$, a
  \emph{replay oracle} is the deterministic counterpart
  \[
    \mathsf{ReplayOracle}\;\mathit{spec}\;\mathit{roles}
    \;\defeq\;
    \mathsf{Counterpart}\;\mathsf{Id}\;\mathit{spec}\;\mathit{roles}\;
      (\lambda\,\_.\;\mathsf{PUnit}).
  \]
  At sender nodes, it observes the sender's move and continues.  At receiver
  nodes, it deterministically chooses a challenge and continues with the
  subtree indexed by that challenge.
  \lean{Interaction.ReplayOracle}
  \uses{int:counterpart}
\end{definition}

Once a replay oracle is fixed, the prover no longer needs to store verifier
challenges explicitly.  It suffices to store only the sender messages; the
receiver moves can be read back from the replay oracle.

\begin{definition}[MessagesOnly and transcript reconstruction]
  \label{int:messages-only}
  Given a replay oracle $\rho : \mathsf{ReplayOracle}\;\mathit{spec}\;
  \mathit{roles}$, the type
  $\mathsf{MessagesOnly}\;\mathit{spec}\;\mathit{roles}\;\rho$ stores exactly
  the sender moves of the protocol:
  \begin{itemize}
    \item At $\mathsf{done}$, it is $\mathsf{PUnit}$.
    \item At a sender node, it stores the chosen message together with the
      recursively stored sender messages of the resulting subtree.
    \item At a receiver node, it stores no new data and recurses directly into
      the subtree determined by $\rho$'s chosen challenge.
  \end{itemize}
  The function
  $\mathsf{MessagesOnly.deriveTranscript}$ reconstructs the full transcript by
  combining the stored sender moves with the replay oracle's receiver
  challenges.
  \lean{Interaction.MessagesOnly}
  \lean{Interaction.MessagesOnly.deriveTranscript}
  \uses{int:replay-oracle, int:transcript}
\end{definition}

This is the point where the dependent interaction tree pays off.  The type of
the tail at a receiver node is not ``the same protocol with one more slot
filled''; it is literally the subtree selected by the replayed challenge.  In
other words, transcript reconstruction is not a post-processing pass over a
flat vector of messages, but a structurally recursive walk through the same
interaction tree as the original protocol.

\subsection{Why verifier-side Fiat--Shamir needs more than a counterpart}

On the prover side, the replay oracle alone is enough.  On the verifier side,
however, we need to \emph{replay} a prescribed transcript through the original
interactive verifier.  This is exactly where the ordinary counterpart type is
too weak.

Recall that an ordinary receiver node in
$\mathsf{Counterpart}\;m\;\mathit{spec}\;\mathit{roles}\;\mathit{Output}$ has
the shape
\[
  m\bigl((x : X) \times \mathit{Cont}(x)\bigr).
\]
This is perfect for execution: it tells us how to sample the next challenge and
continue.  But it does \emph{not} tell us how to recover the continuation for a
\emph{prescribed} challenge $x$ without actually running the sampler.

\begin{definition}[PublicCoinCounterpart]
  \label{int:public-coin-counterpart}
  A \emph{public-coin counterpart} factors each receiver node into two pieces:
  \[
    \mathsf{sample} : m\,X,
    \qquad
    \mathsf{next} : (x : X) \to \mathsf{PublicCoinCounterpart}(\mathit{rest}(x)).
  \]
  Sender nodes are unchanged: the verifier simply observes the sender's move
  and continues.  This factorization captures the public-coin property needed
  for Fiat--Shamir: all verifier randomness is used only to sample the next
  challenge, and the rest of the verifier is a deterministic continuation in
  that challenge.
  \lean{Interaction.Spec.PublicCoinCounterpart}
  \uses{int:counterpart}
\end{definition}

Two generic operations come with this factorization:

\begin{definition}[Forgetting and replaying public-coin structure]
  \label{int:public-coin-replay}
  A public-coin counterpart supports:
  \begin{itemize}
    \item $\mathsf{toCounterpart}$, which forgets the factorization and
      recovers an ordinary executable counterpart by sampling a challenge and
      following its continuation;
    \item $\mathsf{replay}$, which follows a prescribed transcript and ignores
      the samplers entirely, using only the continuation family.
  \end{itemize}
  \lean{Interaction.Spec.PublicCoinCounterpart.toCounterpart}
  \lean{Interaction.Spec.PublicCoinCounterpart.replay}
  \uses{int:public-coin-counterpart}
\end{definition}

This is the missing ingredient that makes verifier-side Fiat--Shamir honest in
the new framework.  We do not transform an arbitrary interactive verifier into
a Fiat--Shamir verifier; we transform a verifier whose receiver nodes are
already known to be public-coin in this replayable sense.

\subsection{Public-coin verifiers and reductions}

\begin{definition}[PublicCoinVerifier and PublicCoinReduction]
  \label{int:public-coin-verifier-reduction}
  A \emph{public-coin verifier} is a statement-indexed family of public-coin
  counterparts.  A \emph{public-coin reduction} is a reduction whose prover is
  unchanged, but whose verifier is public-coin in this sense.
  Forgetful maps
  $\mathsf{PublicCoinVerifier.toVerifier}$ and
  $\mathsf{PublicCoinReduction.toReduction}$ recover the ordinary interaction
  objects.
  \lean{Interaction.PublicCoinVerifier}
  \lean{Interaction.PublicCoinVerifier.toVerifier}
  \lean{Interaction.PublicCoinReduction}
  \lean{Interaction.PublicCoinReduction.toReduction}
  \uses{int:public-coin-replay, int:reduction}
\end{definition}

\subsection{The basic Fiat--Shamir transform}

With these ingredients in place, the transform itself is straightforward.

\begin{definition}[Running the prover against a replay oracle]
  \label{int:run-with-replay-oracle}
  $\mathsf{Strategy.runWithReplayOracle}$ executes a prover strategy against a
  replay oracle.  At sender nodes, it records the prover's chosen move in the
  resulting $\mathsf{MessagesOnly}$ proof.  At receiver nodes, it reads the
  challenge from the replay oracle and continues without storing any additional
  proof data.
  \lean{Interaction.Strategy.runWithReplayOracle}
  \uses{int:messages-only, int:strategy-with-roles}
\end{definition}

\begin{definition}[Interaction-native Fiat--Shamir]
  \label{int:interaction-fiat-shamir}
  The transformed statement is the pair of the original statement with a replay
  oracle.  The transformed protocol is a single sender node carrying a
  $\mathsf{MessagesOnly}$ proof.

  The prover-side transform runs the original prover against the replay oracle:
  \[
    \mathsf{Prover.fiatShamir}.
  \]
  The verifier-side transform takes a public-coin verifier, reconstructs the
  full transcript from the replay oracle and messages-only proof, and replays
  that transcript through the original verifier:
  \[
    \mathsf{PublicCoinVerifier.fiatShamir}.
  \]
  Combining the two yields
  $\mathsf{PublicCoinReduction.fiatShamir}$.
  \lean{Interaction.Prover.fiatShamir}
  \lean{Interaction.PublicCoinVerifier.fiatShamir}
  \lean{Interaction.PublicCoinReduction.fiatShamir}
  \uses{int:run-with-replay-oracle, int:public-coin-verifier-reduction}
\end{definition}

\begin{remark}[What this formalizes, and what it does not]
  This is the \emph{basic} Fiat--Shamir transform at the interaction level.  It
  formalizes the messages-only collapse of a public-coin interactive protocol
  against a fixed replay oracle.  It does \emph{not} yet model the replay oracle
  itself as a random oracle or sponge; that oracle-level formulation belongs in
  a later security layer.
\end{remark}

\subsection{Comparison with the flat oracle-reduction core}

It is instructive to compare this with the older flat, $\Fin$-indexed
formalization from Section~\ref{sec:fiat_shamir}.

\begin{remark}[Flat core versus dependent interaction core]
  \label{int:fiat-shamir-comparison}
  The two formalizations solve related but genuinely different problems.

  In the old flat oracle-reduction core:
  \begin{itemize}
    \item the protocol shape is fixed in advance by a
      $\mathsf{ProtocolSpec}\;n$;
    \item prover messages and verifier challenges are separated globally by
      index sets $\mathsf{MessageIdx}$ and $\mathsf{ChallengeIdx}$;
    \item the Fiat--Shamir oracle is a family indexed by challenge positions;
    \item the verifier already appears as a function of the \emph{full
      transcript}, so verifier-side Fiat--Shamir is obtained simply by
      reconstructing that transcript and calling the original verifier.
  \end{itemize}

  In the new interaction-native core:
  \begin{itemize}
    \item the protocol shape may depend on earlier moves, so the remainder of
      the interaction is a subtree selected by the actual transcript;
    \item a messages-only proof stores precisely the sender moves along the
      realized path, with no global $\Fin$-indexed tuple of all messages;
    \item transcript reconstruction is structurally recursive on the interaction
      tree and therefore automatically follows challenge-dependent subprotocols;
    \item verifier-side Fiat--Shamir requires an explicit public-coin
      factorization of the verifier counterpart, because an ordinary interactive
      verifier is no longer \emph{already} a function on full transcripts.
  \end{itemize}

  So the new formalization is both more general and more explicit.  It is more
  general because it handles genuinely dependent interaction trees.  It is more
  explicit because it isolates the precise replayability property needed on the
  verifier side, rather than baking that property implicitly into the old
  verifier interface.
\end{remark}

\begin{remark}[Position in the literature]
  \label{int:fiat-shamir-literature}
  The standard Fiat--Shamir literature usually starts from a fixed-round
  public-coin protocol.  A representative example is the treatment of
  multi-round Fiat--Shamir by Attema, Fehr, and Kloo{\ss}
  \cite{AttemaFehrKlooss2023}, where the protocol is a $(2\mu+1)$-move
  interactive proof with fixed challenge sets, and the analysis proceeds via
  reconstructed transcript trees.  In that setting, verifier-side
  Fiat--Shamir naturally appears as ``reconstruct the transcript, then run the
  original verifier.''

  Our interaction-native formulation agrees with that picture when the protocol
  has a fixed round schedule, but it identifies an additional structural point
  that becomes unavoidable in the dependent setting: replaying a transcript
  through the verifier requires not just a verifier algorithm, but a
  \emph{replayable public-coin continuation structure}.  This is what
  $\mathsf{PublicCoinCounterpart}$ makes explicit.

  Formal verification work on public-coin proofs, such as the Isabelle
  formalization of sumcheck by Bosshard, Bootle, and Sprenger
  \cite{BBS24}, is closely adjacent in spirit, but it still treats a
  fixed-round protocol rather than a dependent interaction tree.  On the
  engineering side, recent transcript-specification efforts such as the CFRG
  Fiat--Shamir draft \cite{CFRGFiatShamir2025} and the Decree library
  \cite{Decree2024} emphasize staged transcript discipline and challenge-order
  correctness.  Our formulation can be viewed as a semantic counterpart to that
  engineering discipline: in a dependent protocol, requesting the next
  challenge does not merely extend a transcript prefix, but selects the
  continuation subtree in which the rest of the verifier lives.
\end{remark}

\begin{remark}[Why the new formulation is worth keeping]
  The flat formulation remains useful for textbook-style oracle reductions with
  a fixed round schedule, and it is the right home for the random-oracle-model
  security story already developed there.  But the interaction-native
  formulation identifies the structural essence of Fiat--Shamir more cleanly:
  a deterministic replay oracle for transcript reconstruction, and a replayable
  public-coin verifier for the verifier side.  That is the right abstraction
  boundary for protocols whose later message types genuinely depend on earlier
  challenges.
\end{remark}
